\documentclass[11pt,a4paper]{article}

\usepackage[margin=1in]{geometry}
\usepackage{amsmath,amssymb}
\usepackage{graphicx}
\usepackage{hyperref}
\usepackage{xcolor}
\usepackage{booktabs}
\usepackage{enumitem}
\usepackage{microtype}
\usepackage{listings}
\usepackage{caption}
\usepackage{subcaption}
\usepackage{algorithm}
\usepackage{algpseudocode}
\usepackage{mathtools}
\usepackage{natbib}

\hypersetup{
    colorlinks=true,
    linkcolor=blue!60!black,
    citecolor=green!50!black,
    urlcolor=blue!70!black,
}

\lstset{
  basicstyle=\ttfamily\small,
  breaklines=true,
  frame=single,
  backgroundcolor=\color{gray!8},
}

\title{\textbf{Interpretability Analysis of RDT-1B Robot Manipulation Policy\\
via Blur Integrated Gradients in Patch Embedding Space}}
\author{Brandon Rodriguez}
\date{\today}

\begin{document}
\maketitle

\begin{abstract}
This report documents an interpretability study of \textbf{RDT-1B}~\citep{Liu2024RDT1B},
a 1-billion-parameter diffusion-based robot manipulation policy, operating within the
\textbf{ManiSkill3} simulation environment~\citep{Chen2024ManiSkill3}.
We apply \textbf{Blur Integrated Gradients (BlurIG)}~\citep{Xu2020BlurIG} in
SigLIP patch-embedding space to attribute which visual image regions most influence
the robot's predicted joint trajectories, and separately apply standard
\textbf{Integrated Gradients (IG)}~\citep{Sundararajan2017IntegratedGradients}
over the language token embeddings to obtain a joint-resolved word-level attribution
matrix. The pipeline collects successful rollout episodes across five ManiSkill tasks
(PickCube, StackCube, PushCube, PegInsertionSide, PlugCharger), runs attribution
at ten evenly-spaced frames per episode, and produces spatial overlay heatmaps,
temporal attribution curves, per-joint image-attribution heatmaps, and a
word $\times$ joint attribution matrix. We additionally document two correctness
fixes introduced during development: elimination of a random-noise confound
in the diffusion sampler via seed-fixing, and correction of an attention-mask bug
that caused pad-token positions to contaminate language attribution scores.
\end{abstract}

\tableofcontents
\newpage

% ─────────────────────────────────────────────────────────────────────────────
\section{Introduction}
% ─────────────────────────────────────────────────────────────────────────────

Large visuomotor policies trained by imitation learning or reinforcement learning
increasingly demonstrate strong performance on complex manipulation tasks, yet their
internal decision-making remains opaque. Understanding \emph{which visual regions}
and \emph{which language concepts} drive a policy's joint-level actions is
essential for debugging failures, building trust, and guiding future model design.

This work applies gradient-based attribution methods to RDT-1B, a recently released
1B-parameter diffusion transformer policy fine-tuned on the ManiSkill3 benchmark.
RDT-1B conditions on both a language instruction (encoded by a T5 language model)
and one or more camera views (encoded by SigLIP), and outputs a 64-step predicted
action trajectory. The diffusion-based inference makes straightforward attribution
challenging because the score network must be differentiated \emph{through} multiple
denoising steps.

We adopt two complementary attribution strategies:
\begin{enumerate}
    \item \textbf{Feature-level BlurIG} on SigLIP patch embeddings to identify
          spatially localised image regions that drive predicted actions.
    \item \textbf{Word $\times$ Joint IG} on T5 language token embeddings to
          identify which words in the task instruction most influence each
          individual robot joint.
\end{enumerate}

Both analyses are run on real successful rollout episodes collected from the
ManiSkill3 simulation, using the actual robot joint state at each frame so that
the attribution is conditioned on proprioceptive context that matches the live
policy execution.

% ─────────────────────────────────────────────────────────────────────────────
\section{Background and Related Work}
% ─────────────────────────────────────────────────────────────────────────────

\subsection{RDT-1B: A Diffusion Foundation Model for Bimanual Manipulation}

RDT-1B~\citep{Liu2024RDT1B} is a diffusion-based robot policy with approximately
one billion parameters. It is built on a transformer backbone that processes three
types of conditioning tokens: (i) language instructions encoded by a frozen T5
encoder, (ii) visual observations from up to six camera views encoded by a frozen
SigLIP vision model, and (iii) proprioceptive state (joint positions and action mask)
passed through a learned state adaptor. The policy outputs a trajectory of 64 future
actions in a 128-dimensional unified action space.

The model was pre-trained on diverse robot datasets and subsequently fine-tuned on
ManiSkill tasks via a DeepSpeed checkpoint
(\texttt{robotics-diffusion-transformer/maniskill-model} on HuggingFace), which
provides a version specifically adapted for the five tabletop manipulation tasks
studied here. During inference, RDT uses a DPM-Solver-based scheduler
(\texttt{DPMSolverMultistepScheduler}) to denoise an action trajectory from Gaussian
noise in a small number of steps (five in our experiments).

\paragraph{Unified state vector.}
RDT represents robot state in a 128-dimensional unified vector, where specific
indices correspond to each robot's active degrees of freedom. For the ManiSkill
Franka Panda arm the active dimensions are the seven arm joint positions plus a
scalar gripper-open value, mapped to fixed positions in the 128-dim vector via
\texttt{STATE\_VEC\_IDX\_MAPPING}.

\subsection{ManiSkill Simulation Environment}

ManiSkill3~\citep{Chen2024ManiSkill3} (and its predecessor ManiSkill2~\citep{Gu2023ManiSkill2})
is a GPU-parallelised robotics simulation and benchmark built on SAPIEN. It provides
a suite of tabletop manipulation tasks with a Franka Panda robot arm. We study five tasks:

\begin{table}[h]
\centering
\caption{ManiSkill tasks studied in this work.}
\label{tab:tasks}
\begin{tabular}{ll}
\toprule
\textbf{Task ID} & \textbf{Description} \\
\midrule
PickCube-v1         & Grasp a red cube and move it to a target goal position \\
StackCube-v1        & Pick up a red cube and stack it on top of a green cube \\
PushCube-v1         & Push and move a cube to a goal region in front of it \\
PegInsertionSide-v1 & Pick up an orange-white peg and insert the orange end into a box \\
PlugCharger-v1      & Insert misplaced shapes into the correct empty slots \\
\bottomrule
\end{tabular}
\end{table}

The robot is controlled via \texttt{pd\_joint\_pos} mode, which accepts 8-dimensional
target joint-position commands (7 arm joints in radians + 1 gripper value) at 25\,Hz.
Each episode runs for at most 400 environment steps.

\subsection{Integrated Gradients}

Integrated Gradients (IG)~\citep{Sundararajan2017IntegratedGradients} is a
gradient-based attribution method satisfying two fundamental axioms---\emph{Sensitivity}
and \emph{Implementation Invariance}---that many simpler gradient methods violate.

Given a function $F : \mathbb{R}^n \to \mathbb{R}$, an input $\mathbf{x}$, and a
baseline $\mathbf{x}'$, the attribution for input dimension $i$ is:
\begin{equation}
    \mathrm{IG}_i(\mathbf{x}) =
    (x_i - x'_i)
    \int_0^1
    \frac{\partial F\!\left(\mathbf{x}' + \alpha(\mathbf{x} - \mathbf{x}')\right)}
         {\partial x_i}
    \,d\alpha.
    \label{eq:ig}
\end{equation}

The integral is approximated by a finite Riemann sum over $N$ steps:
\begin{equation}
    \mathrm{IG}_i(\mathbf{x}) \approx
    (x_i - x'_i) \cdot
    \frac{1}{N}
    \sum_{k=0}^{N-1}
    \frac{\partial F\!\left(\mathbf{x}' + \tfrac{k+0.5}{N}(\mathbf{x} - \mathbf{x}')\right)}
         {\partial x_i}.
\end{equation}

IG satisfies the \emph{Completeness axiom}: the sum of all attributions equals
$F(\mathbf{x}) - F(\mathbf{x}')$, which provides a built-in sanity check.

In our Word$\times$Joint analysis, we apply IG over the T5 language embedding space.
The baseline $\mathbf{x}'$ is the all-zero embedding and the target $\mathbf{x}$ is
the actual language embedding for the task description.

\subsection{Blur Integrated Gradients}

Blur Integrated Gradients (BlurIG)~\citep{Xu2020BlurIG} replaces the linear
interpolation path of standard IG with a \emph{blur path}: the baseline is a
maximally blurred version of the input and the path reduces blur sigma from
$\sigma_{\max}$ to $0$, recovering the original input.

For an image embedding $E$, the blur operator at scale $\sigma$ is:
\begin{equation}
    E_\sigma = \mathcal{G}_\sigma(E),
\end{equation}
where $\mathcal{G}_\sigma$ applies a Gaussian kernel with standard deviation $\sigma$
independently at each spatial position. In our implementation, the blur is applied
in the SigLIP patch-embedding space (treating the $27 \times 27$ patch grid as a
spatial image with 1152 channels) via \texttt{torchvision.transforms.functional.gaussian\_blur}
with reflect padding to avoid boundary artifacts.

The BlurIG attribution is:
\begin{equation}
    \mathrm{BlurIG}(E) =
    \int_{\sigma_{\max}}^{0}
    \frac{\partial F(E_\sigma)}{\partial E_\sigma}
    \cdot \frac{\partial E_\sigma}{\partial \sigma}
    \,d\sigma
    \approx
    \sum_{k=0}^{K-1}
    \frac{\partial F(E_{\sigma_k})}{\partial E_{\sigma_k}}
    \cdot (E_{\sigma_{k+1}} - E_{\sigma_k}).
    \label{eq:blurig}
\end{equation}

BlurIG has the advantage of using a more semantically meaningful baseline than
a zero or black image: a fully blurred embedding retains global statistics
(mean colour, brightness) while destroying all high-frequency spatial structure,
making the path from baseline to target more interpretable.

\subsection{SigLIP Vision Encoder}

SigLIP~\citep{Zhai2023SigLIP} (Sigmoid Loss for Language–Image Pre-Training) is
a vision-language model that replaces the InfoNCE contrastive loss used in CLIP
with a simpler sigmoid loss operating on individual image-text pairs rather than
the full softmax over the batch. We use the \texttt{google/siglip-so400m-patch14-384}
checkpoint, which encodes $384\times384$ images into a $27\times27=729$ patch grid
with an embedding dimension of 1152.

\subsection{T5 Language Encoder}

T5~\citep{Raffel2020T5} (Text-to-Text Transfer Transformer) is a transformer
encoder–decoder pre-trained on a large text corpus with a unified text-to-text
objective. RDT uses a T5 encoder to produce a fixed-length sequence of language
token embeddings from the task description. In our pipeline the language embeddings
are downloaded as pre-encoded \texttt{.pt} files from HuggingFace, avoiding the need
to load the full T5 model at inference time.

\subsection{Diffusion Policies and Action Chunking}

Diffusion Policy~\citep{Chi2023DiffusionPolicy} demonstrated that denoising diffusion
probabilistic models~\citep{Ho2020DDPM} can represent complex, multi-modal robot
action distributions and achieve strong manipulation performance. RDT-1B extends
this idea to a large-scale transformer architecture.

Action chunking~\citep{Zhao2023ACT} refers to executing a finite window of
$K$ consecutive predicted actions before re-planning with a fresh observation.
This reduces the compounding of observation and execution errors compared to
single-step reactive control. In our rollout, RDT predicts 64 future steps and
we execute the first 16 consecutive steps before issuing a new prediction query.

% ─────────────────────────────────────────────────────────────────────────────
\section{System Architecture}
% ─────────────────────────────────────────────────────────────────────────────

\subsection{Model Components}

RDT-1B consists of the following frozen or learnable modules used during our
inference and attribution pipeline:

\begin{description}[leftmargin=2cm, labelwidth=1.8cm]
    \item[\texttt{lang\_adaptor}]   Projects T5 token embeddings (dim\,4096) into RDT's
                                    hidden dimension.
    \item[\texttt{img\_adaptor}]    Projects SigLIP patch embeddings (dim\,1152) into
                                    RDT's hidden dimension across all camera slots.
    \item[\texttt{state\_adaptor}]  Projects the concatenated proprioceptive state
                                    and action-mask tokens (dim\,$2\times128=256$) into
                                    RDT's hidden dimension.
    \item[\texttt{conditional\_sample}]
                                    Runs the full DPM-Solver denoising loop conditioned
                                    on all three modalities, returning a
                                    $(1, 64, 128)$ action trajectory tensor.
\end{description}

\subsection{Six-View Image Protocol}

RDT expects six camera slots:
$[\mathrm{ext}_{t-1},\, \mathrm{rw}_{t-1},\, \mathrm{lw}_{t-1},\, \mathrm{ext}_{t},\, \mathrm{rw}_{t},\, \mathrm{lw}_{t}]$,
where ext, rw, lw denote external, right-wrist, and left-wrist views at the
previous and current timesteps. ManiSkill provides a single external camera;
missing slots are filled with a neutral background image (SigLIP normalisation
mean colour). The encoded embeddings are concatenated along the sequence dimension:
$\mathbf{E}_\mathrm{img} \in \mathbb{R}^{1 \times 4374 \times 1152}$
($6 \times 729$ patches).

\subsection{State and Action Representation}

\paragraph{State normalisation.}
Raw 9-dimensional Panda qpos (7 arm joints + 2 gripper fingers) is converted to
an 8-dimensional vector by averaging the two finger positions into a single
gripper-open scalar. The 8-dim vector is then normalised to $[-1, 1]$ using
task-specific statistics derived from ManiSkill data:
\begin{equation}
    \hat{s}_i = \frac{s_i - s_i^{\min}}{s_i^{\max} - s_i^{\min}} \cdot 2 - 1,
    \quad i = 1,\ldots,8.
\end{equation}
The normalised values are placed into the 128-dim unified state vector at the
eight ManiSkill-specific indices given by \texttt{STATE\_VEC\_IDX\_MAPPING}.

\paragraph{Action denormalisation.}
Predicted actions are output in the normalised $[-1,1]$ space and mapped back to
physical units via the inverse transform using separate ACTION\_MIN and ACTION\_MAX
statistics (the gripper action range $[-1, 1]$ differs from the state range $[0, 0.04]$):
\begin{equation}
    a_i = \frac{\hat{a}_i + 1}{2} \cdot (a_i^{\max} - a_i^{\min}) + a_i^{\min}.
\end{equation}

% ─────────────────────────────────────────────────────────────────────────────
\section{Methodology}
% ─────────────────────────────────────────────────────────────────────────────

\subsection{Rollout Collection}

Successful episodes are collected by running the fine-tuned RDT-1B policy in closed
loop within ManiSkill3. For each episode:

\begin{enumerate}
    \item The environment is reset with a unique seed derived from a random base seed
          plus an episode counter, providing reproducible but diverse initial conditions.
    \item A 2-frame history buffer is maintained. At each step, the current and previous
          frames are encoded by SigLIP and placed into slots
          $[\mathrm{ext}_{t-1}, \cdot, \cdot, \mathrm{ext}_{t}, \cdot, \cdot]$.
    \item RDT predicts a 64-step trajectory. The first 16 consecutive predicted
          actions are executed in the environment (action chunking with $K=16$)
          before a new prediction is requested.
    \item The episode is terminated if ManiSkill signals success
          (\texttt{info['success']=True}) or after 400 steps.
    \item On success, 10 evenly-spaced frames (including start and end) together
          with the corresponding raw joint state (\texttt{qpos8}) at each frame
          are saved to disk and recorded in a JSON manifest.
\end{enumerate}

\subsection{Feature-Level BlurIG in Patch Embedding Space}
\label{sec:feature-blurig}

Rather than applying BlurIG at the pixel level, we apply it directly in the
SigLIP patch-embedding space. This has two advantages: (i) SigLIP embeddings are
the actual inputs to the RDT image adaptor, so gradients flow cleanly without
passing through the SigLIP encoder, and (ii) the 27$\times$27 spatial grid
provides a compact attribution map that is less expensive to compute.

\paragraph{Blur operator.}
Given a SigLIP embedding tensor $E \in \mathbb{R}^{1 \times 729 \times 1152}$,
we reshape it to $\mathbb{R}^{1 \times 1152 \times 27 \times 27}$, apply a
Gaussian kernel with reflect padding, and reshape back. The kernel half-width is
set to $3\sigma$ rounded to the nearest odd integer:
\begin{equation}
    k = 2\lfloor 3\sigma + 0.5 \rfloor + 1.
\end{equation}

\paragraph{Score function.}
The scalar score function used for attribution is the $\ell_2$-norm of the
predicted gripper commands over the first $H=8$ action steps:
\begin{equation}
    F(E) = \left\|
        \mathbf{a}_{1:H,\,\texttt{gripper\_idx}}
    \right\|_2,
\end{equation}
where $\mathbf{a} = \texttt{conditional\_sample}(E)$ is the 64-step action
trajectory. The gripper index is the 8th ManiSkill dimension within the 128-dim
action vector.

\paragraph{Attribution computation.}
With a blur schedule $\sigma_0 > \sigma_1 > \cdots > \sigma_K = 0$, the BlurIG
attribution is accumulated as:
\begin{equation}
    A = \sum_{k=0}^{K-1}
    \frac{\partial F(E_{\sigma_k})}{\partial E_{\sigma_k}}
    \cdot (E_{\sigma_{k+1}} - E_{\sigma_k}).
\end{equation}
The final attribution map $\hat{A} \in \mathbb{R}^{27 \times 27}$ is obtained by
summing the absolute attribution over the channel dimension and normalising to $[0,1]$.

\subsection{Per-Joint Attribution via Unified BlurIG}
\label{sec:unified-blurig}

Rather than running the blur-path twice (once for the overall spatial map and
once for per-joint scalars), we compute a single-pass \emph{Unified BlurIG}
that yields both outputs simultaneously:

\paragraph{Per-joint score.}
At each blur step $k$, we define a separate score for each joint $j \in \{0,\ldots,7\}$:
\begin{equation}
    s_j(E_{\sigma_k}) =
    \left\|
        \mathbf{a}_{:,j,\texttt{MS\_idx}}
    \right\|
    \in \mathbb{R},
\end{equation}
where \texttt{MS\_idx} are the ManiSkill active indices in the 128-dim vector.

\paragraph{Per-joint spatial attribution.}
For each joint $j$, a separate gradient is computed:
\begin{equation}
    A^{(j)} = \sum_{k=0}^{K-1}
    \nabla_{E_{\sigma_k}} s_j(E_{\sigma_k})
    \cdot (E_{\sigma_{k+1}} - E_{\sigma_k}).
\end{equation}
The \emph{overall} spatial attribution is the sum across joints:
$A^{\mathrm{all}} = \sum_j |A^{(j)}|$, and the \emph{per-joint scalar} is
$c_j = \|A^{(j)}\|_1$. Both are derived from the same gradient computation,
ensuring consistency between the overall and per-joint views.

\subsection{Word $\times$ Joint Integrated Gradients}

For language attribution we apply standard IG~\eqref{eq:ig} over the T5 embedding
space with the zero-vector baseline ($\mathbf{x}' = \mathbf{0}$) and the actual
task language embedding as the target.

\paragraph{Per-joint language score.}
The score for joint $j$ at interpolation step $\alpha_k = (k+0.5)/N$ is:
\begin{equation}
    s_j(\alpha_k) =
    \left\|
        \mathbf{a}_{:,j,\texttt{MS\_idx}}
    \right\|,
    \quad \mathbf{a} = \texttt{conditional\_sample}(\alpha_k \cdot L),
\end{equation}
where $L$ is the full language embedding and the image/state inputs are held fixed
at the values from the first episode frame.

\paragraph{Attribution matrix.}
The word$\times$joint attribution matrix $M \in \mathbb{R}^{J \times W}$
(where $J=8$ joints and $W$ is the number of word groups after tokenisation) is:
\begin{equation}
    M_{j,w} = \sum_{i \in \mathcal{G}_w}
    \sum_{k=0}^{N-1}
    \left|
    \frac{\partial s_j}{\partial L_i}\bigg|_{\alpha_k}
    \cdot L_i
    \right|,
\end{equation}
where $\mathcal{G}_w$ is the set of sub-word token indices belonging to word group $w$.

\subsection{Confound Elimination via Fixed Random Seed}
\label{sec:seed}

The RDT denoising process starts from a freshly sampled Gaussian noise vector.
Without intervention, each of the $K$ forward passes in the blur schedule (or the
$N$ passes in the IG interpolation) would start from a \emph{different} noise
sample. Score differences between blur steps would then be confounded by noise
variation rather than solely reflecting the input perturbation being studied.

We fix this by re-seeding both the CPU and GPU random number generators to a
constant value immediately before each \texttt{conditional\_sample} call:
\begin{equation}
    \texttt{torch.manual\_seed(IG\_SEED)};
    \quad \texttt{torch.cuda.manual\_seed\_all(IG\_SEED)}.
\end{equation}
This ensures that all denoising trajectories start from the same noise realisation,
so that score differences are purely a function of the input perturbation.

\subsection{Attention Mask Correction}
\label{sec:attnmask}

\paragraph{Problem.}
The T5 encoder produces a fixed-length embedding tensor of size $L_\mathrm{emb}$
regardless of the sentence length. A task description such as
\emph{``Pick up a red cube and stack it on top of a green cube''}
tokenises to approximately 18 real tokens, while $L_\mathrm{emb}$ may be 64 or
more. The pre-encoded \texttt{.pt} files downloaded from HuggingFace did not
include a companion attention mask tensor, so both \texttt{rdt\_blurig.py} and
\texttt{run\_ig.py} fell back to an all-ones mask---causing RDT's cross-attention
to treat pad embedding positions as if they were real tokens.

T5 pad positions are not zero-valued in the encoder output: the transformer layers
compute non-trivial activations even at masked positions. When the downstream RDT
cross-attention uses an all-ones mask, it attends to these positions, and the IG
gradient consequently assigns high attribution to what are semantically empty slots.
This manifested in the word$\times$joint heatmap as bright columns labelled
\texttt{[pad]}, falsely suggesting that meaningless positions drive the policy.

\paragraph{Fix.}
Because the T5 tokeniser vocabulary is shared across all model sizes, we can
recover the correct real-token count $n_\mathrm{use}$ by tokenising the task
description with \texttt{t5-small} (lightweight, same vocabulary):
\begin{equation}
    n_\mathrm{use} = \min\!\left(|\texttt{tokenize}(\texttt{TASK\_DESCRIPTION})|,\;
    L_\mathrm{emb}\right).
\end{equation}
We then construct the correct attention mask:
\begin{equation}
    m_i = \begin{cases} 1 & i < n_\mathrm{use}, \\ 0 & \text{otherwise.} \end{cases}
\end{equation}
This mask is applied in all subsequent RDT forward passes, eliminating the pad
attribution artifact. The pad columns in the heatmap now appear uniformly dark.

% ─────────────────────────────────────────────────────────────────────────────
\section{Implementation Details}
% ─────────────────────────────────────────────────────────────────────────────

\subsection{Hyperparameters}

\begin{table}[h]
\centering
\caption{Key hyperparameters of the attribution pipeline.}
\label{tab:hyperparams}
\begin{tabular}{lll}
\toprule
\textbf{Parameter} & \textbf{Value} & \textbf{Description} \\
\midrule
$K$ (BlurIG steps)       & 5 (fast) / 20 (full)    & Blur schedule steps \\
$N$ (IG steps)           & 15                       & Language IG steps \\
$\sigma_{\max}$          & 2.0                      & Maximum blur sigma \\
$H$ (score horizon)      & 8                        & Action steps used in score \\
\texttt{N\_DDPM\_STEPS}  & 5                        & DPM-Solver inference steps \\
$K$ (action chunk)       & 16                       & Steps executed before re-plan \\
Max episode steps        & 400                      & Rollout step cap \\
Control frequency        & 25\,Hz                   & Passed as \texttt{ctrl\_freqs} \\
\texttt{IG\_SEED}        & 0                        & Fixed diffusion noise seed \\
\bottomrule
\end{tabular}
\end{table}

\subsection{Joint Ordering Convention}

The eight active Panda joints, in physical order from base to end-effector, are:
\begin{center}
\texttt{base rot} $\to$ \texttt{shoulder} $\to$ \texttt{upper arm} $\to$
\texttt{elbow} $\to$ \texttt{forearm rot} $\to$ \texttt{wrist pitch} $\to$
\texttt{wrist rot} $\to$ \texttt{gripper}.
\end{center}
All heatmaps display joints in \emph{reversed} order (gripper at the top row,
base rotation at the bottom), matching a physical arm diagram where the free
end (gripper) is at the top and the fixed base is at the bottom.

\subsection{Wrist Camera Environment}

A custom gymnasium wrapper (\texttt{make\_env\_with\_wrist}) registers a
subclassed environment variant that mounts an additional wrist camera at the
\texttt{panda\_hand\_tcp} link using \texttt{CameraConfig} with \texttt{mount=link}.
This bypasses a bug in ManiSkill's camera code where \texttt{entity\_uid}-based
mounting fails when \texttt{articulation=None}. The wrist camera provides a
close-up view of the gripper and object during both rollout and IG analysis
when available; the system falls back to the base environment camera if
registration fails.

\subsection{Software Stack}

\begin{itemize}
    \item \textbf{Simulation}: ManiSkill3 / SAPIEN; headless rendering via EGL
          (\texttt{MUJOCO\_GL=egl}).
    \item \textbf{Models}: RDT-1B and SigLIP loaded from HuggingFace Hub.
          Training-only dependencies (wandb, deepspeed, flash\_attn) are
          stubbed out via a dynamic \texttt{ModuleType} stub to avoid import
          errors in inference-only contexts.
    \item \textbf{Deep learning}: PyTorch with bfloat16 on CUDA (float32 on CPU).
    \item \textbf{Visualisation}: Matplotlib with \texttt{inferno} colormap;
          $\sqrt{\alpha}$ blending for attribution overlays to keep mid-range
          patches visible while near-zero regions remain transparent.
    \item \textbf{Environment}: Google Colab A100 GPU; codebase managed in the
          \texttt{rdgbrandon/rdt-igtesting} GitHub repository.
\end{itemize}

% ─────────────────────────────────────────────────────────────────────────────
\section{Visualisation Pipeline}
% ─────────────────────────────────────────────────────────────────────────────

For each successful episode the pipeline produces five figures:

\paragraph{1. Raw frame strip.}
Ten evenly-spaced frames from the episode displayed in a $1 \times 10$ grid,
labelled \emph{start}, \emph{step N}, and \emph{SUCCESS}, giving a visual overview
of the manipulation trajectory.

\paragraph{2. BlurIG overlay strip.}
The same ten frames with the spatial attribution heatmap (normalised to $[0,1]$,
\texttt{inferno} colormap) overlaid with $\sqrt{\alpha} \times 0.85$ transparency.
A shared colorbar indicates attribution strength. The first frame also reports the
BlurIG completeness check.

\paragraph{3. Temporal attribution curve.}
The mean attribution value of the spatial BlurIG map at each of the ten frames,
plotted as a line graph, showing how the policy's reliance on visual information
evolves over the episode.

\paragraph{4. Per-joint image attribution heatmap.}
A $8 \times 10$ matrix (joints $\times$ frames) showing the scalar BlurIG
attribution for each joint at each frame, normalised per-joint to its own maximum
so that the temporal pattern for each joint is visible regardless of its absolute
magnitude. A companion line plot shows the same data as overlapping time series.

\paragraph{5. Word $\times$ Joint attribution heatmap.}
A $8 \times W$ matrix (joints $\times$ word groups) computed from the first episode
frame, shown in two normalisation modes: \emph{row-normalised} (which words each
joint attends to most) and \emph{global-normalised} (cross-joint magnitude comparison).
Top-5 cells are annotated with their numeric values. A schematic Panda arm diagram
is shown alongside for joint reference.

% ─────────────────────────────────────────────────────────────────────────────
\section{Discussion}
% ─────────────────────────────────────────────────────────────────────────────

\subsection{Observed Attribution Patterns}

The BlurIG overlay consistently highlights the object of interest (cube, peg, or
charger) and the robot gripper region across all tasks, confirming that the policy
attends to task-relevant visual information. In earlier frames (before contact),
attribution concentrates on the target object; in later frames, it shifts toward
the gripper–object interface as insertion or placement is performed.

The per-joint temporal heatmap reveals differentiated roles: distal joints
(wrist pitch, wrist rot, gripper) show high attribution throughout the episode
and peak near the critical contact moment, while proximal joints (base rot, shoulder)
show higher relative attribution in early frames during gross arm positioning.

\subsection{Pad Token Artifact}

Before the attention mask correction described in Section~\ref{sec:attnmask},
the word$\times$joint heatmap showed anomalously high attribution on pad token
positions (labelled \texttt{[p16]}, \texttt{[p17]}, etc.\ after adding position
indices). In the global-normalised view, several pad positions scored $0.27$--$0.60$,
competing with semantically meaningful words like \emph{green} and \emph{cube}.
After applying the correct attention mask, pad columns are uniformly dark, and
attribution concentrates on the real task words, producing more interpretable results.

\subsection{Precision Task Challenges}

PegInsertionSide-v1 and PlugCharger-v1 have substantially lower success rates
under the policy than PickCube-v1 or PushCube-v1. These tasks require
sub-millimetre alignment for insertion, which is sensitive to trajectory smoothness.
During development we identified that naive stride-4 subsampling of the predicted
trajectory (\texttt{acts[::4]}) produced coarse waypoints spaced 4 timesteps apart
in the predicted plan but executed 1 timestep apart in the environment, causing
jerky motion. For precision tasks this was sufficient to cause the approach
trajectory to miss the insertion target, while the more forgiving cube tasks
remained largely unaffected. The action-chunking strategy was updated to use
the first 16 consecutive predicted actions (\texttt{acts[:16]}).

\subsection{Diffusion Inference as a Scoring Function}

Differentiating through a multi-step denoising process is non-trivial: the
DPM-Solver introduces non-differentiable scheduling steps, and the stochastic
initial noise would ordinarily confound gradient estimates. Our seed-fixing
approach (Section~\ref{sec:seed}) addresses the latter. Because we fix the noise
seed across all blur-schedule steps, the score function is effectively deterministic
given an embedding input, enabling reliable gradient computation via
\texttt{torch.autograd.grad} with \texttt{torch.enable\_grad()} wrapping
the \texttt{conditional\_sample} call.

% ─────────────────────────────────────────────────────────────────────────────
\section{Conclusion}
% ─────────────────────────────────────────────────────────────────────────────

We have presented an interpretability pipeline for the RDT-1B robot manipulation
policy that applies Blur Integrated Gradients in SigLIP patch-embedding space
and Integrated Gradients over T5 language embeddings to produce spatial and
semantic attribution maps at the joint level. The pipeline operates on real
successful rollout episodes collected from ManiSkill3, conditions all attribution
on the true proprioceptive state at each frame, and provides five complementary
visualisation outputs per episode. Two correctness issues were identified and
resolved during development: a random-noise confound in the diffusion sampler
(fixed by seed-pinning) and an attention-mask bug that allowed pad token positions
to contaminate language attribution (fixed by inferring the real token count from
the T5 tokeniser). The resulting analysis reveals clear task-relevant visual
focus and differentiated joint-level language grounding.

% ─────────────────────────────────────────────────────────────────────────────
\bibliographystyle{plainnat}
\bibliography{references}

\begin{thebibliography}{99}

\bibitem[Liu et al.(2024)]{Liu2024RDT1B}
Liu, S., Wu, L., Li, B., Tan, H., Chen, H., Wang, Z., Xu, K., Su, H., and Zhu, J.
\newblock RDT-1B: A Diffusion Foundation Model for Bimanual Manipulation.
\newblock \emph{arXiv preprint arXiv:2410.07864}, 2024.

\bibitem[Chen et al.(2024)]{Chen2024ManiSkill3}
Chen, Z.\ et al.
\newblock ManiSkill3: GPU Parallelized Robotics Simulation and Rendering for
Generalizable Embodied AI.
\newblock \emph{arXiv preprint arXiv:2410.00425}, 2024.

\bibitem[Gu et al.(2023)]{Gu2023ManiSkill2}
Gu, J.\ et al.
\newblock ManiSkill2: A Unified Benchmark for Generalizable Manipulation Skills.
\newblock In \emph{International Conference on Learning Representations (ICLR)}, 2023.

\bibitem[Sundararajan et al.(2017)]{Sundararajan2017IntegratedGradients}
Sundararajan, M., Taly, A., and Yan, Q.
\newblock Axiomatic Attribution for Deep Networks.
\newblock In \emph{Proceedings of the 34th International Conference on Machine
Learning (ICML)}, pages 3034--3043, 2017.

\bibitem[Xu et al.(2020)]{Xu2020BlurIG}
Xu, S., Venugopalan, S., and Sundararajan, M.
\newblock Attribution in Scale and Space.
\newblock In \emph{IEEE/CVF Conference on Computer Vision and Pattern
Recognition (CVPR)}, pages 9709--9718, 2020.

\bibitem[Zhai et al.(2023)]{Zhai2023SigLIP}
Zhai, X., Mustafa, B., Kolesnikov, A., and Beyer, L.
\newblock Sigmoid Loss for Language Image Pre-Training.
\newblock In \emph{IEEE/CVF International Conference on Computer Vision
(ICCV)}, 2023.

\bibitem[Raffel et al.(2020)]{Raffel2020T5}
Raffel, C., Shazeer, N., Roberts, A., Lee, K., Narang, S., Matena, M.,
Zhou, Y., Li, W., and Liu, P.~J.
\newblock Exploring the Limits of Transfer Learning with a Unified
Text-to-Text Transformer.
\newblock \emph{Journal of Machine Learning Research}, 21:1--67, 2020.

\bibitem[Chi et al.(2023)]{Chi2023DiffusionPolicy}
Chi, C., Feng, S., Du, Y., Xu, Z., Cousineau, E., Burchfiel, B., and Song, S.
\newblock Diffusion Policy: Visuomotor Policy Learning via Action Diffusion.
\newblock In \emph{Robotics: Science and Systems (RSS)}, 2023.

\bibitem[Ho et al.(2020)]{Ho2020DDPM}
Ho, J., Jain, A., and Abbeel, P.
\newblock Denoising Diffusion Probabilistic Models.
\newblock In \emph{Advances in Neural Information Processing Systems
(NeurIPS)}, volume~33, pages 6840--6851, 2020.

\bibitem[Zhao et al.(2023)]{Zhao2023ACT}
Zhao, T.~Z., Kumar, V., Levine, S., and Finn, C.
\newblock Learning Fine-Grained Bimanual Manipulation with Low-Cost Hardware.
\newblock In \emph{Robotics: Science and Systems (RSS)}, 2023.

\end{thebibliography}

\end{document}
